{\small\textbf{Risks addressed (Diego, 1:15) [12:25–13:40]}}\par\smallskip
So which financial-stability risks was the buffer meant to address? DNB's reports list them, and each risk was matched to a tool.
\begin{itemize}\setlength\itemsep{0pt}
\item \textbf{Shocks nobody can forecast} go to the countercyclical buffer, because it can be released whatever the source: the war, an energy shock, a sudden tightening of financial conditions, debt-sustainability problems, or bank turmoil abroad like SVB and Credit Suisse in 2023.
\item \textbf{The known hot spot, housing}, went to other tools: borrower-based limits, \textcolor{navy}{\textbf{[def]}} caps on loan-to-value and loan-to-income, tax reform, and a floor on mortgage risk weights.
\item \textbf{The size of the big banks} goes to the buffer for systemically important banks.
\end{itemize}
\par\smallskip
\textcolor{navy}{\textbf{[def]}} This is the Tinbergen principle: one instrument per target. DNB itself said that more releasable capital is valuable "given the sensitivity of the Dutch economy to external events".
